\documentclass[12pt]{article}
\input{format}

\usepackage{titlesec}

% --- Main Section Formatting ---
\titleformat{\section}
  {\normalfont\Large\bfseries}{\thesection}{1em}{}
\titlespacing*{\section}{0pt}{2\baselineskip}{1\baselineskip}

% --- Subsection Formatting (NEW) ---
\titleformat{\subsection}
  {\normalfont\large\bfseries} % \large is smaller than \Large
  {\thesubsection}             
  {1em}                        
  {}                           

\titlespacing*{\subsection}
  {0pt}               % Left margin
  {1.5\baselineskip}  % Space BEFORE the subsection (1.5 lines)
  {0.5\baselineskip}  % Space AFTER the subsection (half a line)

\begin{document}
\title{SLAM - Ex 5}
\author{Iris Kaplan \& William Ayoub}
\maketitle

This document contains the answers to VAN Exercise $5$. The code for this exercise is in \href{https://github.com/iriskaplan/SLAM_Project}{this link}. The final version used for this submission corresponds to commit
\href{todo: add!!11}
{\texttt{}}.



\section{Track Reprojection Error}

In this section we evaluate one feature track using the camera poses estimated in Exercise 3. The global coordinate system was chosen as the coordinate system of the first left camera. The relative PnP transformations from Exercise 3 were composed into global extrinsic camera matrices $T_{0 \to i}$, where $i$ is the frame index. These matrices were saved by \texttt{section\_3\_6} in \texttt{code/ex3.py}, line $304$.

We loaded the tracking database constructed in Exercise 4 and randomly selected a track of length $15$. For every frame in this track, we used the corresponding global camera matrix to construct a \texttt{gtsam.StereoCamera}.

Let the extrinsic camera matrix of a frame be $[R \mid t]$, where $R \in \mathbb{R}^{3 \times 3}$ and $t \in \mathbb{R}^{3}$. In our convention, this matrix maps points from the global coordinate system to the camera coordinate system:
$$
x_{\text{cam}} = R x_{\text{world}} + t.
$$
However, \texttt{gtsam.StereoCamera} expects the pose in the opposite direction, from camera coordinates to global coordinates. Since $R^{-1}=R^T$, we get
$$
x_{\text{world}}
=
R^{-1}(x_{\text{cam}} - t)
=
R^T x_{\text{cam}} - R^T t.
$$
Therefore, the pose passed to GTSAM is
$$
\left[
R^T \mid -R^T t
\right].
$$

Using the \texttt{StereoCamera} of the last frame in the selected track, we triangulated a 3D point with \texttt{backproject}. Since the camera pose is given in the camera-to-global convention, the triangulated point is represented in global coordinates. We then projected this point into all frames of the track using \texttt{project}, and compared the projections with the measured stereo locations from the tracking database.

For frame $i$, let the measured and projected stereo points be
$$
z_i =
\begin{bmatrix}
u_{L,i} \\
u_{R,i} \\
v_i
\end{bmatrix},
\qquad
\hat z_i =
\begin{bmatrix}
\hat u_{L,i} \\
\hat u_{R,i} \\
\hat v_i
\end{bmatrix}.
$$
The stereo reprojection error plotted in Figure~\ref{fig:reprojection} is
$$
\|z_i - \hat z_i\|_2.
$$

\begin{figure}[H]
    \centering
    \includegraphics[width=1\linewidth]{william_plots/q5_1_reprojection_error.png}
    \caption{Stereo reprojection error over the selected track. The 3D point was triangulated from the last frame, so the error at that frame is approximately zero. The figure was generated by \texttt{question5\_1} in \texttt{code/ex5.py}, line 65, using \texttt{plot\_reprojection\_errors\_q5\_1} in \texttt{code/slam/visualization/ex5\_plots.py}, lines 21--50.}
    \label{fig:reprojection}
\end{figure}

We also created one \texttt{gtsam.GenericStereoFactor3D} for each frame projection. We used an isotropic pixel noise model for the stereo coordinates $(u_L,u_R,v)$, with standard deviation $\sigma = 1$ pixel. Thus, the covariance matrix was
$$
\Sigma = \sigma^2 I_3 = I_3.
$$

Denoting the residual by $r_i = z_i - \hat z_i$, the stereo factor error is
$$
E_i = \frac{1}{2} r_i^T \Sigma^{-1} r_i.
$$
Since $\Sigma = I_3$, this becomes
$$
E_i = \frac{1}{2}\|r_i\|_2^2.
$$
Thus, the factor error is one half of the squared reprojection error. This relation is visible in Figures~\ref{fig:reprojection} and~\ref{fig:factor}, which have the same qualitative shape.

\begin{figure}[H]
    \centering
    \includegraphics[width=1\linewidth]{william_plots/q5_1_factor_error.png}
    \caption{GTSAM stereo factor error over the selected track. The figure was generated by \texttt{question5\_1} in \texttt{code/ex5.py}, line 65, using \texttt{plot\_factor\_errors\_q5\_1} in \texttt{code/slam/visualization/ex5\_plots.py}, lines 53--82.}
    \label{fig:factor}
\end{figure}

\section{Transformation in a Common Coordinate System}

Let cameras $A$ and $B$ have extrinsic matrices $[R_A \mid t_A]$ and $[R_B \mid t_B]$, respectively, with both extrinsics defined with respect to the same global coordinate system. Thus, for a global point $x_g$, its coordinates in camera $A$ are
$$
x_A = R_A x_g + t_A,
$$
and its coordinates in camera $B$ are
$$
x_B = R_B x_g + t_B.
$$

We want to express the transformation from the coordinate system of camera $A$ to the coordinate system of camera $B$. First, we invert the transformation of camera $A$:
$$
x_g = R_A^T(x_A - t_A).
$$

Substituting this expression into the equation for camera $B$ gives
$$
x_B
=
R_B R_A^T(x_A - t_A) + t_B.
$$
Rearranging terms, we obtain
$$
x_B
=
R_B R_A^T x_A
+
\left(t_B - R_B R_A^T t_A\right).
$$

Therefore, the transformation (and extrinsic matrix) from camera $A$ coordinates to camera $B$ coordinates is
$$
T_{A \to B}
=
\left[
R_B R_A^T
\mid
t_B - R_B R_A^T t_A
\right].
$$

\section{Bundle Adjustment Window}

For simplicity, we selected keyframes at a fixed interval of $10$ frames, resulting in local windows of $11$ consecutive stereo pairs.

The first bundle window spans frames $0$--$10$. It contains $11$ camera poses, $1810$ landmarks, and $7710$ factors, including one prior factor anchoring the first camera pose. Table~\ref{tab:ba_first_window_stats} summarizes the factor-graph error before and after optimization.

\begin{table}[h]
    \centering
    \caption{Optimization statistics for the first bundle window. The average factor error is the total graph error divided by the total number of factors. Values are printed by \texttt{\_print\_bundle\_summary()} in \texttt{code/ex5.py}, lines 97--109.}
    \label{tab:ba_first_window_stats}
    \begin{tabular}{lc}
        \hline
        Quantity & Value \\
        \hline
        Initial total graph error & $8433.007$ \\
        Final total graph error & $1551.122$ \\
        Number of factors & $7710$ \\
        Initial average factor error & $1.094$ \\
        Final average factor error & $0.201$ \\
        \hline
    \end{tabular}
\end{table}

The optimization reduces the total graph error by approximately a factor of
$5.4$.

We next inspected the projection factor with the largest initial error. This
factor connects camera pose $c=3$ with landmark $q=861$.

\begin{table}[h]
    \centering
    \caption{Largest initial-error stereo projection factor before and after
    bundle adjustment.}
    \label{tab:largest_factor_diagnostic}
    \begin{tabular}{lcc}
        \hline
        Quantity & Initial values & Optimized values \\
        \hline
        Factor error & $244.491$ & $1.468$ \\
        Left-image distance & $15.412$ px & $1.713$ px \\
        Right-image distance & $17.395$ px & $1.663$ px \\
        \hline
    \end{tabular}
\end{table}

For both the initial and optimized estimates, we constructed a
\texttt{StereoCamera} using the corresponding pose of $c$ and projected the
corresponding landmark $q$. Figure~\ref{fig:largest_reproj_err} compares the
measured feature position with the projected position in the left and right
images.

\begin{figure}[h]
    \centering
    \includegraphics[width=1\linewidth]{william_plots/q5_3_largest_initial_factor.png}
    \caption{Stereo reprojection diagnostic for the factor with the largest
    initial error. The red cross denotes the stereo measurement and the green
    circle denotes the projected landmark. The top row uses the initial values
    and the bottom row uses the optimized values. The figure was generated by
    \texttt{question5\_3} in \texttt{code/ex5.py}, line 191, using
    \texttt{plot\_largest\_factor\_diagnostic} in
    \texttt{code/slam/visualization/ex5\_plots.py}, lines 117--200.}
    \label{fig:largest_reproj_err}
\end{figure}

Figures~\ref{fig:ba_first_window_top_down} and~\ref{fig:ba_first_window_3d}
show the optimized scene for the first bundle window. Camera poses and landmarks
are expressed relative to frame $0$, which is the first frame of this specific
window.

\begin{figure}[H]
    \centering
    \includegraphics[width=1\linewidth]{william_plots/bundle_scene_top_down.png}
    \caption{Top-down view of the optimized first bundle window. The red
    trajectory shows the optimized camera centers and the blue points show
    optimized landmarks. Both axes are measured in meters and use equal scale.
    The figure was generated by \texttt{question5\_3} in \texttt{code/ex5.py},
    line 204, using \texttt{plot\_bundle\_scene\_top\_down} in
    \texttt{code/slam/visualization/trajectory.py}, lines 102--139.}
    \label{fig:ba_first_window_top_down}
\end{figure}

\begin{figure}[H]
    \centering
    \includegraphics[width=1\linewidth]{william_plots/bundle_scene_3d.png}
    \caption{Three-dimensional view of the optimized first bundle window,
    including the camera trajectory and reconstructed landmarks. Axes are
    measured in meters. The figure was generated by \texttt{question5\_3} in
    \texttt{code/ex5.py}, line 203, using \texttt{plot\_bundle\_scene\_3d} in
    \texttt{code/slam/visualization/trajectory.py}, lines 52--99.}
    \label{fig:ba_first_window_3d}
\end{figure}

\clearpage
\section{Sliding Bundle Windows}

We chose all keyframes along the trajectory at a fixed interval of $10$ frames and solved all resulting bundle windows independently. Each window uses the coordinate system of its first frame. The initial values for poses and landmarks were transformed from global coordinates into the first-frame-local coordinate system. A total of $330$ bundle windows were solved.

\subsection{First Frame Position of the Last Bundle Window}
The last bundle window spans frames $3290$--$3299$. After optimization, the exact position of the first frame of this window is $[-2.838 \cdot 10^{-14}, -7.081 \cdot 10^{-15}, -7.306 \cdot 10^{-15}]^T$, which is essentially $[0, 0, 0]^T$ up to machine precision.

\subsection{Anchoring Factor Final Error}
The anchoring factor's final error is exactly $0.0$. This is expected because the prior factor strongly constrains the first pose to its initial value (which is defined as the identity in the local coordinate system of the bundle window), and the optimizer can satisfy this constraint completely with zero cost.

\subsection{View from Above with Keyframes and 3D Points}
From each solved bundle window we extracted the relative poses and composed them into absolute global poses. Figure~\ref{fig:all_keyframes_top_down} shows the view from above (2D) of the scene, presenting all left-camera keyframes (in blue) and the optimized 3D landmarks (light blue points). 

\subsection{Overlay with Ground Truth Poses}
As requested, Figure~\ref{fig:all_keyframes_top_down} also overlays the estimated keyframes with the ground-truth keyframe poses from the KITTI dataset (shown in orange). 

\begin{figure}[H]
    \centering
    \includegraphics[width=1\linewidth]{william_plots/q5_4_keyframes_and_landmarks_top_down.png}
    \caption{Top-down view of all bundle-adjusted keyframes (blue) and ground-truth keyframe positions (orange). Optimized landmarks are shown as light blue scatter points. The figure was generated by \texttt{question5\_4} in \texttt{code/ex5.py}, line 277, using \texttt{plot\_keyframes\_and\_landmarks\_top\_down} in \texttt{code/slam/visualization/ex5\_plots.py}, lines 203--261.}
    \label{fig:all_keyframes_top_down}
\end{figure}

\subsection{Keyframe Localization Error}
Figure~\ref{fig:localization_error} presents the keyframe localization error in meters (location difference only, Euclidean distance) over time. The error is defined as $e_k = \left\| \hat{p}_k - p_k^{\text{GT}} \right\|_2$.

The mean keyframe localization error is $22.159$~m. The error grows and oscillates over the sequence, reflecting drift accumulation in the composed relative poses since we perform only local bundle adjustment without explicit loop closure.

\begin{figure}[H]
    \centering
    \includegraphics[width=1\linewidth]{william_plots/q5_4_keyframe_localization_error.png}
    \caption{Keyframe localization error (Euclidean distance between estimated and ground-truth positions) over the frame index. The figure was generated by \texttt{question5\_4} in \texttt{code/ex5.py}, line 285, using \texttt{plot\_keyframe\_localization\_error} in \texttt{code/slam/visualization/ex5\_plots.py}, lines 264--292.}
    \label{fig:localization_error}
\end{figure}

\end{document}
